\begin{explanationbox}
	\textbf{\textcolor{CExplanation}{一句话直觉}}
	不进行角度或边长计算，仅通过向量的数量积与模长运算，即可判定三角形是直角、等腰还是等边。

	\medskip
	\textbf{\textcolor{CExplanation}{核心拆解}}
	\begin{description}[style=nextline, leftmargin=2.8em, labelsep=0.5em, itemsep=0.45em]
		\item[\textbf{要点一（数量积定直角）：}]
			数量积为零 $\Longleftrightarrow$ 两向量垂直。因此若 $\overrightarrow{AB}\cdot\overrightarrow{AC}=0$，则 $AB\perp AC$，即 $\angle A=90^\circ$，三角形为直角三角形。
		\item[\textbf{要点二（模长定等腰）：}]
			向量模长相等 $\Longleftrightarrow$ 对应边长相等。若 $|\overrightarrow{AB}|=|\overrightarrow{AC}|$，则 $AB=AC$，三角形为等腰三角形。
		\item[\textbf{要点三（模长全等定等边）：}]
			三个模长全相等 $\Longleftrightarrow$ 三边全相等。若 $|\overrightarrow{AB}|=|\overrightarrow{AC}|=|\overrightarrow{BC}|$，则三角形为等边三角形。
	\end{description}

	\medskip
	\textbf{\textcolor{CExplanation}{几何本质（向量几何关系）}}
	数量积 $\vec{a}\cdot\vec{b}=|\vec{a}||\vec{b}|\cos\theta$ 为零时 $\cos\theta=0$，即两向量垂直，对应直角。向量的模 $|\vec{a}|$ 表示线段长度，模长相等等价于边长相等。因此向量运算直接刻画了三角形的边角属性。

	\medskip
	\textbf{\textcolor{CExplanation}{代数意义（坐标运算简化）}}
	在坐标系中，设 $A(x_1,y_1),B(x_2,y_2),C(x_3,y_3)$，则 $\overrightarrow{AB}=(x_2-x_1,y_2-y_1)$，数量积与模长均可通过坐标的加、减、乘、开方快速算出。整个判定过程完全代数化。

	\medskip
	\textbf{\textcolor{CExplanation}{考点价值（快速判定形状）}}
	\begin{itemize}[leftmargin=*, itemsep=0.5em, topsep=0.4em]
		\item \textbf{考法一（直接坐标判定）：} 给出三角形顶点坐标，直接计算数量积和模长，判断形状。
		\item \textbf{考法二（条件转化判定）：} 给出向量关系式（如 $\overrightarrow{AB}\cdot\overrightarrow{AC}=0$ 或 $|\overrightarrow{AB}|=|\overrightarrow{AC}|$），直接推出形状。
		\item \textbf{考法三（综合推理判定）：} 在向量综合题中，先利用数量积或模长判定三角形形状，再进一步计算面积、夹角等。
	\end{itemize}

	\medskip
	\textbf{\textcolor{CExplanation}{顿悟点}}
	\textbf{垂直与相等是两类独立的判定}：数量积只处理垂直，模长只处理相等，二者互不干扰。遇到具体条件时只需对号入座，无需同时验证多个条件。

	\medskip
	\textbf{\textcolor{CExplanation}{使用场景}}
	\begin{itemize}[leftmargin=*, itemsep=0.5em, topsep=0.4em]
		\item \textbf{场景一（坐标显式）：} 已知三角形三顶点坐标，求向量后直接计算。
		\item \textbf{场景二（向量等式）：} 已知 $\overrightarrow{AB}\cdot\overrightarrow{AC}=0$ 或 $|\overrightarrow{AB}|=|\overrightarrow{AC}|$ 等条件，立即判定形状。
		\item \textbf{场景三（综合问题）：} 在向量与其他知识综合的问题中，先利用向量关系判定形状再进行后续运算。
	\end{itemize}

\end{explanationbox}
